\chapter{Tại Sao Bạn Bè Làm Em Mệt}
\markboth{Tại Sao Bạn Bè Làm Em Mệt}{Why Do Friends Exhaust Me}

\section{Ở Cùng Nhóm Nhưng Lúc Nào Cũng Phải Gồng}
\begin{truyen}{Ở Cùng Nhóm Nhưng Lúc Nào Cũng Phải Gồng}{Being in a Group but Always Performing}
\chuhoa{C}{ó những học sinh đi với nhóm bạn rất đông nhưng lúc nào cũng phải cười đúng lúc, nói đúng nhịp, theo đúng gu của nhóm.}
\textit{(Some students stay in large friend groups but constantly feel the need to laugh on cue, speak in the right rhythm, and match the group's style.)}

Ở gần người khác mà lúc nào cũng phải gồng thì rất mệt.
\textit{(Being near others while always performing is exhausting.)}

Có em nói một câu nghe rất người lớn: “Ở với bạn mà em còn mệt hơn ở một mình.”
\textit{(One student said something strangely mature: “Being with friends tires me out more than being alone.”)}

Tình bạn không cho phép mình là mình sẽ dần trở thành nơi tiêu hao.
\textit{(Friendship that does not allow authenticity slowly becomes draining.)}

Có em không bỏ nhóm vì sợ còn mệt hơn nếu bị ra ngoài.
\textit{(Some students do not leave the group because exclusion feels even scarier.)}

\end{truyen}

\begin{baihoc}
Không phải tình bạn đông người nào cũng là chỗ nghỉ cho tâm hồn.
\textit{(Not every large friendship circle is a resting place for the heart.)}
\end{baihoc}

\begin{ghinhoanh}
\textbf{Short English to remember:}

Friendship becomes tiring when you cannot be yourself inside it.

\end{ghinhoanh}

\begin{tuhocnhanh}
\textbf{Gợi ý để dạy và tự nghĩ / Teaching and reflection prompts}

1. Vì sao ở trong nhóm mà vẫn mệt? \textit{Why can students feel exhausted even inside a group?}

2. Điều gì khiến một tình bạn thành nơi tiêu hao? \textit{What makes a friendship draining?}

3. Mình có đang giúp học sinh nhận ra điều này không? \textit{Am I helping students recognize this?}
\end{tuhocnhanh}
\ngancach

\section{Bị Đùa Hoài Ở Một Điểm Yếu}
\begin{truyen}{Bị Đùa Hoài Ở Một Điểm Yếu}{Being Repeatedly Teased at the Same Weak Spot}
\chuhoa{C}{ó những câu trêu lặp đi lặp lại về ngoại hình, giọng nói, nhà cửa, hay chuyện học.}
\textit{(Some jokes repeat again and again about appearance, voice, family, or school performance.)}

Ban đầu người bị trêu còn cười.
\textit{(At first the target may still laugh.)}

Sau đó là mệt. Rồi tránh. Rồi thu lại.
\textit{(Later comes exhaustion, then avoidance, then withdrawal.)}

Bạn bè có thể không đánh nhau mà vẫn làm nhau đau rất đều.
\textit{(Friends may never fight openly and still hurt each other steadily.)}

\end{truyen}

\begin{baihoc}
Một chỗ yếu bị chạm lặp lại sẽ thành nơi một học sinh mất dần bình an.
\textit{(A weak spot touched repeatedly becomes a place where a student slowly loses peace.)}
\end{baihoc}

\begin{ghinhoanh}
\textbf{Short English to remember:}

Repeated teasing turns a weak spot into a wound.

\end{ghinhoanh}

\begin{tuhocnhanh}
\textbf{Gợi ý để dạy và tự nghĩ / Teaching and reflection prompts}

1. Vì sao đùa lặp lại nguy hiểm hơn tưởng tượng? \textit{Why is repeated teasing more harmful than people think?}

2. Học sinh thường phản ứng theo những giai đoạn nào? \textit{What stages do students often go through?}

3. Mình có đang xem nhẹ những trò đùa lặp lại không? \textit{Am I underestimating repeated jokes?}
\end{tuhocnhanh}
\ngancach

\section{Sợ Mất Bạn Nên Phải Làm Theo}
\begin{truyen}{Sợ Mất Bạn Nên Phải Làm Theo}{Fearing the Loss of Friends, So Following Along}
\chuhoa{C}{ó em biết việc nhóm bạn làm là dở, nhưng vẫn theo.}
\textit{(Some students know their friends are doing the wrong thing, but still go along.)}

Không phải vì không biết đúng sai.
\textit{(It is not because they do not know right from wrong.)}

Mà vì mất bạn ở tuổi này nghe còn đáng sợ hơn bị phạt.
\textit{(It is because losing friends at this age feels scarier than punishment.)}

Nỗi sợ bị bỏ ra ngoài làm nhiều học sinh yếu đi ở chỗ vốn tưởng mình sẽ mạnh.
\textit{(The fear of exclusion weakens many students exactly where they believed they would be strong.)}

\end{truyen}

\begin{baihoc}
Muốn dạy học sinh biết nói không, phải hiểu các em đang sợ mất điều gì khi nói không.
\textit{(If we want students to say no, we must understand what they fear losing by saying it.)}
\end{baihoc}

\begin{ghinhoanh}
\textbf{Short English to remember:}

To teach students to say no, first understand what they fear losing.

\end{ghinhoanh}

\begin{tuhocnhanh}
\textbf{Gợi ý để dạy và tự nghĩ / Teaching and reflection prompts}

1. Vì sao học sinh biết sai vẫn làm theo? \textit{Why do students follow even when they know it is wrong?}

2. Các em sợ mất gì nhất? \textit{What are they most afraid of losing?}

3. Mình có đang dạy đủ về áp lực đồng trang lứa không? \textit{Am I teaching enough about peer pressure?}
\end{tuhocnhanh}
\ngancach

\section{Em Muốn Có Bạn Nhưng Em Cũng Muốn Yên}
\begin{truyen}{Em Muốn Có Bạn Nhưng Em Cũng Muốn Yên}{I Want Friends but I Also Want Peace}
\chuhoa{C}{ó những học sinh không cô độc hoàn toàn, nhưng cứ thân với ai một thời gian là mệt.}
\textit{(Some students are not completely lonely, but every friendship starts to feel exhausting after a while.)}

Các em cần gần gũi, nhưng cũng cần không gian.
\textit{(They need closeness, but they also need space.)}

Tuổi này chưa phải em nào cũng biết nói ranh giới của mình cho tử tế.
\textit{(At this age, not every student knows how to express boundaries well.)}

Nhiều tình bạn hỏng không phải vì hết quý, mà vì không biết giữ khoảng cách cho vừa.
\textit{(Many friendships break not because affection is gone, but because healthy distance was never learned.)}

\end{truyen}

\begin{baihoc}
Tình bạn bền không chỉ nhờ gần nhau, mà còn nhờ biết chừa chỗ cho nhau thở.
\textit{(Lasting friendship depends not only on closeness, but also on leaving each other room to breathe.)}
\end{baihoc}

\begin{ghinhoanh}
\textbf{Short English to remember:}

Good friendship leaves room to breathe.

\end{ghinhoanh}

\begin{tuhocnhanh}
\textbf{Gợi ý để dạy và tự nghĩ / Teaching and reflection prompts}

1. Vì sao thân thiết cũng có thể gây mệt? \textit{Why can closeness also become tiring?}

2. Học sinh cần học gì về ranh giới? \textit{What should students learn about boundaries?}

3. Mình có giúp các em hiểu vừa gần vừa yên là có thể không? \textit{Am I helping students see that closeness and peace can coexist?}
\end{tuhocnhanh}
\ngancach

\section{Em Cười Theo Đám Bạn Nhưng Về Nhà Rất Mệt}
\begin{truyen}{Em Cười Theo Đám Bạn Nhưng Về Nhà Rất Mệt}{I Laugh with Friends but Go Home Drained}
\chuhoa{C}{ó học sinh trong lớp nhìn rất hòa đồng. Giờ ra chơi nói nhiều, gặp bạn là cười.}
\textit{(Some students look very social at school. They talk a lot at break time and laugh with friends easily.)}

Nhưng về nhà là im ngay.
\textit{(But once home, they fall silent immediately.)}

Tôi hỏi: “Hôm nay vui không?”
\textit{(I asked, “Was today fun?”)}

Em đáp: “Vui thì có, mà mệt nữa.”
\textit{(The student replied, “Fun, yes. But exhausting too.”)}

\end{truyen}

\begin{baihoc}
Có những tình bạn đông vui nhưng vẫn làm học sinh hao rất nhiều sức bên trong.
\textit{(Some lively friendships still consume a great deal of inner energy.)}
\end{baihoc}

\begin{ghinhoanh}
\textbf{Short English to remember:}

Friendship can look lively and still be draining.

\end{ghinhoanh}

\begin{tuhocnhanh}
\textbf{Gợi ý để dạy và tự nghĩ / Teaching and reflection prompts}

1. Vì sao ở trong đám đông vẫn có thể mệt? \textit{Why can students feel drained even in a lively group?}

2. Học sinh đang phải gồng ở phần nào? \textit{What part of themselves are they straining?}

3. Mình có đang nhầm vui vẻ với khỏe mạnh không? \textit{Am I mistaking sociability for emotional health?}
\end{tuhocnhanh}
\ngancach

\section{Bạn Thân Nhưng Cứ Đòi Em Phải Chọn Phe}
\begin{truyen}{Bạn Thân Nhưng Cứ Đòi Em Phải Chọn Phe}{Close Friends Keep Forcing Me to Choose Sides}
\chuhoa{C}{ó nhóm bạn rất thân, nhưng hễ có chuyện là bắt nhau chọn phe.}
\textit{(Some friend groups are close, but every conflict turns into a demand to choose sides.)}

“Nếu chơi với nó nữa thì đừng chơi với tao.”
\textit{(“If you keep hanging out with her, then don't hang out with me.”)}

Câu kiểu đó làm một đứa trẻ mệt khủng khiếp.
\textit{(That kind of line is exhausting for a child.)}

Nhiều em còn chưa biết giữ mình, đã phải giữ cả quan hệ của mấy người cùng lúc.
\textit{(Many students barely know how to protect themselves, yet are forced to manage several relationships at once.)}

\end{truyen}

\begin{baihoc}
Tình bạn tốt không ép người khác phản bội người thứ ba để chứng minh sự trung thành.
\textit{(Good friendship does not force betrayal to prove loyalty.)}
\end{baihoc}

\begin{ghinhoanh}
\textbf{Short English to remember:}

Good friendship does not demand betrayal as proof of loyalty.

\end{ghinhoanh}

\begin{tuhocnhanh}
\textbf{Gợi ý để dạy và tự nghĩ / Teaching and reflection prompts}

1. Vì sao việc bắt chọn phe làm học sinh kiệt? \textit{Why does forced side-taking exhaust students?}

2. Dấu hiệu của tình bạn kiểm soát là gì? \textit{What are signs of controlling friendship?}

3. Mình có dạy học sinh nhận ra ranh giới này không? \textit{Am I teaching students to recognize this boundary?}
\end{tuhocnhanh}
\ngancach

\section{Em Sợ Mất Bạn Nên Không Dám Nói Thật}
\begin{truyen}{Em Sợ Mất Bạn Nên Không Dám Nói Thật}{I Fear Losing Friends, So I Do Not Tell the Truth}
\chuhoa{C}{ó em bị bạn mượn đồ, mượn tiền, nói quá lời, nhưng vẫn cười cho qua.}
\textit{(Some students let friends borrow things, money, or cross lines, and still laugh it off.)}

Tôi hỏi: “Sao em không nói thẳng?”
\textit{(I asked, “Why don't you say it directly?”)}

Em đáp: “Tại em sợ nói ra rồi không còn chơi với ai.”
\textit{(The student answered, “Because I'm afraid if I say it, I won't have anyone left.”)}

Rất nhiều sự nhịn của tuổi học trò không đến từ rộng lượng, mà từ sợ cô độc.
\textit{(Much of school-age tolerance comes not from generosity, but from fear of loneliness.)}

\end{truyen}

\begin{baihoc}
Muốn dạy học sinh nói thật, phải hiểu các em đang sợ trả giá điều gì cho sự thật đó.
\textit{(If we want students to speak honestly, we must understand what price they fear paying for it.)}
\end{baihoc}

\begin{ghinhoanh}
\textbf{Short English to remember:}

Fear of loneliness often keeps students silent.

\end{ghinhoanh}

\begin{tuhocnhanh}
\textbf{Gợi ý để dạy và tự nghĩ / Teaching and reflection prompts}

1. Vì sao học sinh chọn nhịn thay vì nói thật? \textit{Why do students choose silence over honesty?}

2. Nỗi sợ nào đang đứng sau sự nhịn này? \textit{What fear stands behind this silence?}

3. Mình có thể giúp các em nói thật an toàn hơn bằng cách nào? \textit{How can I help them speak honestly more safely?}
\end{tuhocnhanh}
\ngancach

\section{Có Bạn Chơi Chung Chứ Không Có Bạn Đứng Về Em}
\begin{truyen}{Có Bạn Chơi Chung Chứ Không Có Bạn Đứng Về Em}{I Have People to Hang Out With, but No One to Stand by Me}
\chuhoa{C}{ó học sinh không hẳn cô đơn. Vẫn có người ăn cùng, nói chuyện cùng, đi cùng.}
\textit{(Some students are not exactly alone. They still have people to eat, talk, and walk with.)}

Nhưng đến lúc bị oan, bị trêu quá mức, hay bị đẩy ra ngoài, không ai đứng về.
\textit{(But when they are treated unfairly, mocked too much, or pushed out, no one stands beside them.)}

Khi đó các em mới hiểu sự khác nhau giữa chơi chung và thật sự có bạn.
\textit{(That is when they learn the difference between hanging out and truly having a friend.)}

Sự hụt này làm tuổi học trò lạnh đi rất nhanh.
\textit{(That gap can turn school life cold very quickly.)}

\end{truyen}

\begin{baihoc}
Tình bạn thật được lộ ra rõ nhất không phải lúc vui, mà lúc mình bị đẩy về phía yếu thế.
\textit{(Real friendship shows itself most clearly not in fun times, but when we are pushed into a weaker position.)}
\end{baihoc}

\begin{ghinhoanh}
\textbf{Short English to remember:}

Real friendship appears most clearly when you are vulnerable.

\end{ghinhoanh}

\begin{tuhocnhanh}
\textbf{Gợi ý để dạy và tự nghĩ / Teaching and reflection prompts}

1. Khác nhau giữa chơi chung và có bạn thật là gì? \textit{What is the difference between hanging out and real friendship?}

2. Vì sao khoảnh khắc yếu thế lại lộ bạn thật? \textit{Why does vulnerability reveal true friendship?}

3. Mình có giúp học sinh nhận ra điều này không? \textit{Am I helping students see this difference?}
\end{tuhocnhanh}
\ngancach

\section{Em Muốn Nghỉ Chơi Một Mình Một Hôm Có Được Không}
\begin{truyen}{Em Muốn Nghỉ Chơi Một Mình Một Hôm Có Được Không}{May I Spend One Break Alone}
\chuhoa{C}{ó em chỉ muốn một hôm không nhập nhóm nào, không phải cười, không phải cập nhật chuyện gì hết.}
\textit{(Some students simply want one day without joining any group, smiling on cue, or staying updated on everything.)}

Nhưng tuổi này ở một mình cũng dễ bị hỏi: “Bị gì vậy?”
\textit{(But at this age, even being alone invites questions like, “What's wrong with you?”)}

Có những học sinh rất cần một chút riêng mà vẫn sợ bị hiểu là có vấn đề.
\textit{(Some students badly need a little solitude yet fear being seen as abnormal.)}

Ngay cả trong tình bạn, con người cũng cần chỗ để thở một mình.
\textit{(Even inside friendship, a person still needs room to breathe alone.)}

\end{truyen}

\begin{baihoc}
Tình bạn lành không bắt học sinh phải có mặt đủ mọi lúc để chứng minh mình còn thân.
\textit{(Healthy friendship does not force constant presence as proof of closeness.)}
\end{baihoc}

\begin{ghinhoanh}
\textbf{Short English to remember:}

Healthy friendship leaves room for solitude.

\end{ghinhoanh}

\begin{tuhocnhanh}
\textbf{Gợi ý để dạy và tự nghĩ / Teaching and reflection prompts}

1. Vì sao học sinh ngại xin một khoảng riêng? \textit{Why are students afraid to ask for space?}

2. Tình bạn lành cho phép điều gì? \textit{What does healthy friendship allow?}

3. Mình có đang xem một mình là bất thường không? \textit{Am I treating solitude as something abnormal?}
\end{tuhocnhanh}
\ngancach

\section{Một Câu “Kệ Nó Đi” Có Khi Là Điều Cứu Em}
\begin{truyen}{Một Câu “Kệ Nó Đi” Có Khi Là Điều Cứu Em}{Sometimes “Let It Go” from the Right Friend Saves Me}
\chuhoa{C}{ó khi cả đám đang đẩy một chuyện đi quá xa, chỉ cần một người bạn nói: “Kệ nó đi, đừng ép nữa.”}
\textit{(Sometimes a group is taking something too far, and one friend says, “Leave it alone, don't push anymore.”)}

Một câu ngắn vậy mà chặn được rất nhiều mệt mỏi.
\textit{(One short sentence like that can stop a great deal of exhaustion.)}

Tuổi học trò rất cần những bạn bè biết dừng lại đúng lúc.
\textit{(School age deeply needs friends who know when to stop.)}

Không phải bạn tốt nào cũng làm điều to. Nhiều khi bạn tốt chỉ là người không hùa theo lúc cả đám đang quá tay.
\textit{(Not every good friend does something grand. Sometimes a good friend is simply the one who refuses to join the crowd when it goes too far.)}

\end{truyen}

\begin{baihoc}
Một người biết dừng đúng lúc có thể giữ được phẩm giá cho bạn mình nhiều hơn mình tưởng.
\textit{(A person who stops at the right moment can preserve a friend's dignity more than they realize.)}
\end{baihoc}

\begin{ghinhoanh}
\textbf{Short English to remember:}

A good friend sometimes protects you by refusing to join the crowd.

\end{ghinhoanh}

\begin{tuhocnhanh}
\textbf{Gợi ý để dạy và tự nghĩ / Teaching and reflection prompts}

1. Vì sao vai trò của một người biết dừng lại lại quan trọng? \textit{Why does one person who knows when to stop matter so much?}

2. Hành động nhỏ nào có thể cứu bạn mình khỏi xấu hổ? \textit{What small action can save a friend from humiliation?}

3. Mình có đang dạy học sinh can đảm kiểu này không? \textit{Am I teaching this kind of courage?}
\end{tuhocnhanh}
